% The honest parties as numbered algorithms. Each macro is a ruled float
% followed by the paragraph that walks its lines, as in Compact Shielded
% CSV. main.tex and the short sections call them where they belong: the
% full version places all of them after the predicates (section on the
% abstract protocol, and the receive section), the short version only
% \algExit, \algAccept, \algReceive and \algSettle.
% Source: spec/layertree.md 9.1 to 9.7 (rule numbers in the prose),
% design/graft.md, design/batching.md, design/pending-receive.md,
% design/towers.md. Where the spec is silent the prose says so.

\newcommand{\kw}[1]{\text{\bfseries #1}}
\newcommand{\fn}[1]{\ensuremath{\mathsf{#1}}}
\newcommand{\Dl}{\ensuremath{\Delta(q)}}

% ------------------------------------------------------------------
\newcommand{\algCreate}{%
\begin{algorithm}
\caption{Create \textendash\ operator, founding a tree}\label{alg:create}
\pseudocode[linenumbering,codesize=\scriptsize]{%
\kw{choose}\ D,\ \mathsf{OPERATOR\_KEY} \text{ and the option set } \mathcal{O} \text{ (heights, sizes, variants)} \\
\kw{choose}\ \text{the outpoint } o_0 \text{ that input 0 of the funding transaction will spend} \\
\mathsf{tree\_id} \gets H_{\mathsf{tree\text{-}id}}(o_0) \\
m \gets \text{Merkle root of the taptree for } (D, \mathsf{OPERATOR\_KEY}, \mathsf{tree\_id}, \mathcal{O}) \\
\kw{build}\ \text{the initial leaves, at least one, and their root } r_0 \\
\kw{publish}\ (D, \mathsf{OPERATOR\_KEY}, \mathsf{tree\_id}, \mathcal{O}, m, r_0) \text{ and each member's leaf and path} \\
\kw{after}\ \text{every member has run } \fn{AcceptLeaf}\text{:} \\
\pcind \text{spend } o_0 \text{ at input 0; output 0} \gets (\TreeSpk(r_0, m),\ \textstyle\sum r_0)
}
\end{algorithm}%
Lines 2 to 4 fix the order that the circularity of spec~6.2 forces: the
scripts of \tr{CLAIM} carry $\mathsf{tree\_id}$, so $m$ depends on it, and
the tree output's own outpoint cannot be its source, since the output's
script depends on $m$. The recommended derivation (line~3) hashes the
outpoint spent by input~0 of the funding transaction, which exists before
the transaction is built and can be spent once, so it is unique by
construction; the spec requires only that it be unique per lineage.
Line~6 publishes the parameters from which any party recomputes $m$
(spec~9.1: a participant MUST recompute $m$ before treating an output as a
tree, which also catches a wrong $\mathsf{tree\_id}$). Line~7 funds the
output only after the founding members have verified their paths
(spec~6.5, 9.2), so each founder runs \fn{AcceptLeaf} before any funds
move. The spec leaves the transport of leaves and paths to founders
open (spec~9, introduction), and says nothing of a timeout.
}

% ------------------------------------------------------------------
\newcommand{\algGraft}{%
\begin{algorithm}
\caption{Graft \textendash\ operator, adding a cohort at height $h$}\label{alg:graft}
\pseudocode[linenumbering,codesize=\scriptsize]{%
\kw{pick}\ \text{an empty position } (h, i) \text{ with path } \pi,\ \text{ at root } r \text{ and sum } S \\
\kw{build}\ \text{the subtree of the new leaves; } \bar n = (\bar h, \bar s) \gets \text{its node} \\
r' \gets \fold(\bar n, \pi) \\
\mathsf{tx}\text{: input 0 spends the tree output; output 0} = (S + \bar s, \TreeSpk(r', m)) \\
\kw{add}\ \text{funding inputs (the operator's, or each member's own)} \\
\kw{if}\ \text{more than one funder:}\ \text{add a change output or a P2A anchor} \\
\kw{for each}\ \text{member } U\text{:}\ \kw{send}\ (\mathsf{tx},\ \ell_U,\ \text{path of } \ell_U \text{ in the subtree},\ \pi,\ r,\ m) \\
\kw{collect}\ \text{the signatures of the members' inputs, if any} \\
\kw{if}\ \text{a member withdraws or does not sign:}\ \text{drop it, rebuild, go to line 2} \\
\text{sign input 0 with } \mathsf{OPERATOR\_KEY};\ \kw{broadcast}
}
\end{algorithm}%
Lines 1 to 4 build the transaction that $\Phi_{\tr{JOIN}}$ accepts: the
slot is empty, output~0 carries $S + \bar s$ and the operator signs. The
slot's height fixes the script, since each height $h$ is a tapleaf fixed
at creation (spec~6.2), and the operator chooses the slot that later
batching depends on (design/graft.md). Line~5 covers the three funding
modes: in an operator-funded graft the operator supplies $\bar s$ and is
paid off-chain, so there is nothing to collect and no round (spec~9.2); in a
member-funded one each member adds an input. Line~6 is the spec's advice
for more than one funder: because each funder signs with
\textsf{SIGHASH\_DEFAULT} or \textsf{SIGHASH\_ALL}, the graft cannot be
replaced at a higher fee without every funder signing again, so the
operator must be able to bump it by CPFP (spec~9.2). Line~7 sends each
member what \fn{AcceptLeaf} needs, and lines 8 and 9 are the abort and
retry of a signing round: the output commits to the whole subtree, so one
member who changes its leaf or never signs invalidates every other
signature (spec~9.2, informative). The spec states the round but gives
the operator no rule for when to give up on a member or how long to wait.
}

% ------------------------------------------------------------------
\newcommand{\algAccept}{%
\begin{algorithm}
\caption{AcceptLeaf \textendash\ joiner $U$, before funding or paying for a leaf}\label{alg:accept}
\pseudocode[linenumbering,codesize=\scriptsize]{%
\kw{input:}\ \text{unsigned } \mathsf{tx},\ \ell_U,\ \text{subtree path } \rho,\ \text{path } \pi \text{ of the slot},\ r,\ m \\
\kw{if}\ \vin_0 \text{ is not an unspent output with } \spk = \TreeSpk(r, m) \text{ and amount } S\text{:}\ \kw{reject} \\
\kw{if}\ \fold(\Emp_h, \pi) \ne (r, S)\text{:}\ \kw{reject} \\
\kw{if}\ |\rho| \ne h \text{ or } \fold(\nd{\ell_U}, \rho) \ne (\bar h, \bar s) \text{ (all sums recomputed)}\text{:}\ \kw{reject} \\
r' \gets \fold((\bar h, \bar s), \pi) \\
\kw{if}\ \vout_0 \ne (S + \bar s, \TreeSpk(r', m))\text{:}\ \kw{reject} \\
\kw{if}\ \ell_U \text{'s fields or } U \text{'s inputs and change differ from the agreed, or} \\
\pcind (\ell_U.\ok, \ell_U.\nc, \ell_U.\sq) \text{ was held in this lineage before}\text{:}\ \kw{reject} \\
\kw{if}\ U \text{ has an input to sign:}\ \text{sign it } (\mathsf{SIGHASH\_DEFAULT});\ \text{store } \ell_U, \rho, \pi, r, r', \mathsf{txid} \\
\kw{else}\ (\text{pays off-chain}):\ \text{wait for the confirmed graft; } \kw{check}\ \vout_0 \text{ is unspent} \\
\pcind \text{and } \fold(\nd{\ell_U}, \text{its path in } r') = (r', S')\text{;}\ \text{only then pay}
}
\end{algorithm}%
\iffullversion
Lines 2 to 6 are checks (a) to (d) of spec~9.2, and they are the checks the
covenant does not make: $\Phi_{\tr{JOIN}}$ sees a hash and a sum and
cannot tell a cohort from random bytes. Line~2 shows that input~0 is the
unspent tree output, line~3 that the slot was empty in the current root,
line~4 that $U$'s own leaf folds to the grafted node, with the subtree
path of exactly $h$ levels so that the leaf lands at depth $D$ (a leaf
grafted shallower sits where no $D$-level path can reach it), and lines 5
and 6 that output~0 commits to the root containing it. This is the check that Theorem~\ref{thm:ground} and the
claim ``an exit is guaranteed only to a member who verified its path''
rest on. Lines 7 and 8 add the agreed fields and, from spec~9.2, the
rule against re-joining a lineage with an owner key, nonce and $\sq$ held
before, which would revive every exit authorisation ever published for
that leaf. A member who funds an input is then protected by
construction, since the only transaction that spends its input commits to
a root containing its leaf (line~9). A member who pays off-chain has
nothing to sign, so lines 2 to 6 protect nothing and the spec requires
the check of lines 10 and 11 on the \emph{confirmed} transaction before payment;
paying earlier trusts the operator to graft. The alternative that the
covenant enforces is a hash-locked voucher against a leaf the operator
already holds (spec~9.2, Algorithm~\ref{alg:swap}).
\else
Lines 2 to 6 are checks (a) to (d) of spec~9.2, which $\Phi_{\tr{JOIN}}$
cannot make since it sees only a hash and a sum: the output is unspent
(2), the slot was empty (3), $U$'s leaf folds to the grafted node along a
subtree path of exactly $h$ levels (4), and output~0 commits to the root
containing it (5, 6). Lines 7 and 8 add the agreed fields and the rule
against re-joining a lineage with an owner key, nonce and $\sq$ held
before, which would revive every exit authorisation published for that
leaf. A member that funds an input is protected by construction (line~9);
one that pays off-chain has nothing to sign, so lines 10 and 11 must hold
on the \emph{confirmed} graft before it pays.
\fi
}

% ------------------------------------------------------------------
\newcommand{\algExit}{%
\begin{algorithm}
\caption{Exit \textendash\ member $U$ with wallet $(\ell, \pi, r, o)$, or a tower holding $\tau$}\label{alg:exit}
\pseudocode[linenumbering,codesize=\scriptsize]{%
\kw{loop}\ \text{(rebuild)}: \\
\pcind (r', S') \gets \fold(\Emp_0, \pi) \\
\pcind \mathsf{tx}\text{: } \vin_0 = o,\ \text{a fee input};\ \vout_0 = (S - \ell.\amt, \TreeSpk(r', m)) \\
\pcind \vout_1 \gets (v, \Ptr{\ell.\xk}),\ v = \ell.\amt \text{ (less, for a plain exit that pays the fee)} \\
\pcind \sigma \gets \fn{Sign}_{\ell.\ok}(\mathsf{tx}),\ \text{ or } \tau \text{ under } \tr{EXIT}^{t}_K\text{; } \kw{broadcast} \\
\pcind \kw{wait}\ \text{until } \mathsf{tx} \text{ is buried or a spend } \mathsf{tx}' \text{ of } o \text{ confirms:} \\
\pcind\pcind \kw{if}\ \mathsf{tx} \text{ is pending and too cheap:}\ \text{bump by replacement or an anchor child} \\
\pcind\pcind \kw{if}\ \mathsf{tx}' \text{ confirms:}\ (\ell, \pi, r, o) \gets \fn{Follow}(\mathsf{tx}') \\
\pcind\pcind\pcind \kw{if}\ \ell \text{ is gone:}\ \kw{stop};\ \kw{else}\ \kw{continue loop}
}
\end{algorithm}%
\iffullversion
Lines 2 to 4 are $\Phi_{\tr{EXIT}}$ for $K = 1$: the leaf is removed along its own
path (check~2), output~1 pays the exit key (check~3) and output~0
recreates the tree output (check~7); line~4 is the payout rule, in
which a plain exit may pay less than the balance and the exact variants
may not. Line~5 signs all outputs, which is what stops a relayer from
diverting the shortfall or the fee slack (spec~7.2 rule~3). The
signature commits to the tree outpoint $o$, which every transition
changes, so it cannot be made in advance: a plain exit needs the owner
online to rebuild it (lines 1 to 5 again). Under $\tr{EXIT}^{t}$ the owner makes
$\tau = \fn{Sign}_{\ell.\ok}(H_{\mathsf{exit}}(\mathsf{tree\_id}\,\|\,\hat\ell))$ once and hands $(\ell, \tau)$ to a tower, which then runs the whole
algorithm and pays the fee from its own input (spec~7.2 rule~7, 9.3); the
tower also tracks the leaf's bytes, not only its path, since a
\tr{REWRITE} top-up changes the amount the exit must pay without the
owner (spec~9.7). Lines 6 to 8 reflect that every transition spends the
one tree output, so an exit can be overtaken by any conflicting spend
that confirms first and has to be rebuilt against the new root
(spec~10.5), after \fn{Follow} (Algorithm~\ref{alg:follow}) has
brought the path up to date; under~[C] this takes at most $\Dl$ blocks. The spec names
replacement (BIP-125) and anchor outputs as the ways to bump a fee but
does not fix a policy, and for a cooperative exit it advises asking the
operator first, who can fold it into its next transition (spec~9.3). An
owner who stays in the tree after an exit built on a published
authorisation was outbid must change the leaf to void it (spec~7.2
rule~7).
\else
Lines 2 to 4 build $\Phi_{\tr{EXIT}}$ for $K = 1$, and line~5 signs all
outputs, so that a relayer cannot divert the shortfall or the fee slack.
The signature commits to the tree outpoint, which every transition
changes, so a plain exit needs the owner online to rebuild it; under
$\tr{EXIT}^{t}$ the owner signs $\tau$ once and a tower runs the algorithm
with its own fee input (spec~7.2 rule~7, 9.3). Lines 6 to 8 reflect that a
conflicting spend that confirms first forces a rebuild against the new
root (spec~10.5), within $\Dl$ under~[C]; \fn{Follow} updates the path
from that spend's witness alone (spec~9.7). The spec fixes no fee-bumping
policy and advises asking the operator for a cooperative exit first.
\fi
}

% ------------------------------------------------------------------
\newcommand{\algFollow}{%
\begin{algorithm}
\caption{Follow \textendash\ member or tower, on each spend of the tree output}\label{alg:follow}
\pseudocode[linenumbering,codesize=\scriptsize]{%
\kw{input:}\ (\ell, \pi, r, o),\ \text{and a transaction } \mathsf{tx}' \text{ that spends } o \text{ at input 0} \\
\text{read from } \mathsf{tx}'\text{'s witness: changed positions, paths used, old and new nodes} \\
r' \gets \text{the root committed by } \vout_0 \text{ of } \mathsf{tx}' \text{ (} \Emp_D \text{ if the tree is gone)} \\
\kw{if}\ \text{the position of } \ell \text{ changed:}\ \ell \gets \text{the new leaf the witness shows (or none)} \\
\kw{else}\ \text{replace the one entry of } \pi \text{ where the two paths diverge} \\
\pcind \text{by the node the witness computes at that height} \\
\kw{if}\ \ell \ne \text{none and } \fold(\nd{\ell}, \pi) \ne r'\text{:} \\
\pcind \text{ask the operator for the path and } \kw{check}\ \text{it against } r' \\
(r, o) \gets (r', (\mathsf{txid}(\mathsf{tx}'), 0))
}
\end{algorithm}%
Line~2 uses the fact that every transition publishes the sibling paths
it used and its old and new nodes (spec~9.7). Line~5 is the whole of path
maintenance: of the siblings along a member's path, only the ancestor of
the changed position at the level where the two paths diverge can change,
and the witness computes its new value (Lemma~\ref{lem:path}, the Lean
\lean{path\_maintenance}). Line~4 is the case where the leaf itself
changes, which the member's own signature accounts for except where the
amount moves without the owner: an amount-only top-up under \tr{REWRITE},
a merge (spec~7.5.4) or a transfer of the reclaim form (spec~7.9 rule~8),
each of which changes the leaf's amount and nothing else, so a tower that holds an exit authorisation must
track the leaf's bytes (spec~7.9 rule~3, 9.7). Line~7 is the fall-back
the spec allows to a participant that has fallen behind: it can always
catch up from the chain, or obtain the current path from the operator and
verify it against the confirmed root. A light client cannot find the
spends by scriptPubKey, since $\TreeSpk(r, m)$ changes with every
transition; it follows the outpoint chain (design/towers.md). The spec
does not say how often a member must run the algorithm, only that it
MUST process every spend (spec~9.7).
}

% ------------------------------------------------------------------
\newcommand{\algVouch}{%
\begin{algorithm}
\caption{Vouch \textendash\ issuer $I$ of a voucher against its leaf $\ell_s$}\label{alg:vouch}
\pseudocode[linenumbering,codesize=\scriptsize]{%
\kw{on}\ \text{request for a voucher paying } \ell_c \text{ (for } \tr{CLAIM}^{h}\text{, with a payment hash } \eta\text{)}: \\
\pcind \kw{if}\ \text{another leaf has the body of } \ell_s \text{ (dense profile)}\text{:}\ \kw{refuse} \\
\pcind \kw{if}\ \textstyle\sum_{\text{outstanding}} + \ell_c.\amt > \ell_s.\amt\text{:}\ \kw{refuse} \\
\pcind \nu \gets \fn{Sign}_{\ell_s.\ok}\bigl(H_{\mathsf{voucher}}(\mathsf{tree\_id}\,\|\,\hat\ell_s\,\|\,\ell_c [\,\|\,\eta])\bigr) \\
\pcind \kw{record}\ (\nu, \ell_s.\sq) \text{ as a liability};\ \kw{send}\ (\mathsf{tree\_id}, \ell_s, \ell_c, \nu) \\
\kw{on}\ \text{a transition that changes a field of } \ell_s \text{ other than } \amt\text{:} \\
\pcind \text{re-sign every outstanding voucher over the new } \ell_s \text{ and the new amount} \\
\kw{on}\ \text{a decrease of } \ell_s.\amt \text{ that leaves } \sq \text{ alone:} \\
\pcind \text{keep the outstanding sum within it, or advance } \sq \text{ and re-sign} \\
\kw{retire}\ \text{a liability once } \ell_s.\sq \text{ exceeds its value at issuance}
}
\end{algorithm}%
Lines 3 and 5 are the issuer's MUSTs of spec~9.6: never issue vouchers
whose amounts sum to more than the source holds, and track each as a
liability. The covenant prevents over-\emph{redemption} but not
over-issuance (spec~10.3). Line~2 is the bookkeeping rule of the replay
paragraph of spec~7.5: a voucher binds its source by content, not by
position, so two leaves with one body would let one voucher redeem
against both, and the dense profile does not check this on chain. Line~4
binds the voucher to $\mathsf{tree\_id}$, the source without its amount
and the claim leaf in full, and for $\tr{CLAIM}^{h}$ to the payment hash
under a tag of its own, so that a plain and a hash-locked voucher are
disjoint instruments. Lines 6 to 9 follow from what the message binds. A
\tr{CLAIM} or \tr{SPLIT} advances $\sq$, an \tr{UPDATE} or
\tr{REWRITE} may rotate a key, renew a lease or change $\sq$, and each of
these voids every outstanding voucher, so the issuer re-signs them
against the new leaf, which also lowers the balance they draw on. A
change of the amount alone voids nothing: a top-up is harmless, and
after a decrease the issuer must keep the outstanding sum within the
reduced amount or advance $\sq$ and rebuild. Line~10 is how a voucher is
retired, since the covenant gives it no deadline (spec~7.5.1). The spec
does not tell the issuer to check that $\ell_c$ obeys the carve rule or
the dust floor; a voucher that violates either simply fails at
redemption.
}

% ------------------------------------------------------------------
\newcommand{\algRedeem}{%
\begin{algorithm}
\caption{Redeem \textendash\ recipient of a voucher $(\mathsf{tree\_id}, \ell_s, \ell_c, \nu)$}\label{alg:redeem}
\pseudocode[linenumbering,codesize=\scriptsize]{%
\kw{on}\ \text{receipt:}\ \kw{if}\ \nu \text{ does not verify under } \ell_s.\ok \text{ over the voucher message, or} \\
\pcind \mathsf{tree\_id} \text{ is not the value from which } m \text{ was recomputed}\text{:}\ \kw{reject} \\
\kw{obtain}\ \text{a path for a leaf with the body of } \ell_s \text{ in the confirmed tree;} \\
\pcind \kw{if}\ \text{it fails or } \ell_s.\amt < \ell_c.\amt\text{:}\ \kw{reject} \\
\kw{read}\ \ell_s.\ek,\ \ell_s.\ea\text{;}\ \kw{if}\ \ell_s.\ea \text{ can pass before a redemption confirms } (\Dl) \text{ and} \\
\pcind \ell_s.\ek \notin \{\mathsf{NUMS}, U\} \text{ and } \ell_s.\xk \ne U\text{:}\ \text{treat as unsecured credit} \\
\kw{if}\ \ell_c.\amt < \ell_s.\amt \text{ and no empty position is known in the source's cohort:} \\
\pcind \text{prefer a full-balance voucher} \\
\kw{to redeem}\text{:}\ \kw{build}\ \mathsf{tx} \text{ for } \Phi_{\tr{CLAIM}} \text{ (position, } \nu, \text{ fee input, preimage if any);} \\
\pcind \kw{broadcast};\ \text{rebuild as in } \fn{Exit} \text{ on a conflicting spend; settled once it confirms}
}
\end{algorithm}%
Lines 1 to 4 are the recipient's MUSTs of spec~9.6: verify the signature,
check that $\mathsf{tree\_id}$ is one of the published creation
parameters from which $m$ is recomputed (spec~9.1), and verify that a
leaf with the voucher's source body is in the tree with an amount that
covers the voucher. Lines 5 and 6 are the lease check. Eviction voids a
voucher without the issuer doing anything, since the evictor can pay the
whole balance to the source's exit key once $\ell_s.\ea$ has passed, so a
recipient treats a voucher whose lease can pass within $\Dl$ as unsecured
credit and should want the lease to outlast its own redemption deadline
plus $\Dl$; the check is moot when the evictor is $\mathsf{NUMS}$ or the
recipient, or when the exit key is the recipient's, the credit leaf of
spec~10.3. Lines 7 and 8 are the empty-position caveat: a voucher for less than
the whole balance needs an empty position for its claim leaf, which the
operator controls. Lines 9 and 10 are the redemption, which only the voucher's
signature authorises; it can be built by anyone who holds the voucher,
a tower included. The spec states that the voucher is settled only after
the redemption, or after the issuer includes the claim leaf and the
recipient verifies inclusion, and that until then the recipient bears the
risk that the issuer double-issues (spec~9.6). For a hash-locked voucher
the recipient does not treat it as payment until it also holds the
preimage and redeems promptly, because the issuer can void it by
advancing $\sq$ (spec~9.6.1).
}

% ------------------------------------------------------------------
\newcommand{\algSwap}{%
\begin{algorithm}
\caption{Swap \textendash\ issuer $A$ and holder $B$ of a $\tr{CLAIM}^{h}$ voucher paired with an HTLC (an ordered exchange)}\label{alg:swap}
\pseudocode[linenumbering,codesize=\scriptsize]{%
A\text{:}\ \text{draw } R;\ \eta \gets \mathsf{SHA256}(R);\ \text{issue } B \text{ a } \tr{CLAIM}^{h} \text{ voucher and an invoice on } \eta \\
B\text{:}\ \kw{if}\ \text{the source's lease can pass before } \mathsf{cltv\_expiry} + \Dl\text{:}\ \kw{abort} \\
B\text{:}\ \kw{pay}\ \text{the invoice (offer the HTLC)} \\
A\text{:}\ \kw{from}\ \text{the offer until the voucher is redeemed or the HTLC fails or times out:} \\
\pcind \text{do not advance } \ell_s.\sq \\
A\text{:}\ \text{settle the HTLC off-chain, releasing } R \text{ through the payment} \\
B\text{:}\ \kw{on}\ \text{learning } R\text{:}\ \text{redeem the voucher at once (} \fn{Redeem}\text{)} \\
B\text{:}\ \text{call the swap complete only when the redemption is buried so deep} \\
\pcind \text{that no reorganisation can replace it by a transition advancing } \sq \\
A\text{:}\ \kw{if}\ \text{the HTLC failed or timed out:}\ \text{retire the voucher by advancing } \ell_s.\sq
}
\end{algorithm}%
The hash orders the two legs and does not make the exchange atomic: the
voucher is a cheque, the value stays in $\ell_s$, and until the
redemption confirms the issuer can void it by a conflicting spend of that
leaf (spec~9.6.2). The five rules of the spec fix the order. Line~1:
the party the HTLC pays draws the preimage (rule~1), because it reveals
it to collect the HTLC and the revelation enables the voucher; were the
voucher holder to draw it, it could redeem and never settle. Lines 6 and 7:
the voucher is the second leg to resolve (rule~2), since it has no
deadline and the HTLC does, and an HTLC must never wait on a preimage
that only an on-chain redemption would reveal; where that order cannot
be arranged the parties use two HTLCs. Lines 4 and 5 are rule~3, which no
script enforces and which binds the issuer even after the HTLC settles:
an issuer that collects and then voids has been paid for a voucher that
can no longer be redeemed. Line~2 is rule~5, the lease check before the
holder offers the HTLC. Lines 7 to 9 are rule~4: the holder carries the
issuer risk until its redemption confirms and is buried, and the
interval includes contention for the tree output, so it is governed by
$\Dl$ and not by a constant. Line~10 retires the voucher as spec~9.6.1
requires, and the issuer must not reuse $\eta$ against the same leaf
until it has. Usually only the voucher reaches the chain, at the cost of
one $\tr{CLAIM}^{h}$.
}

% ------------------------------------------------------------------
\newcommand{\algRewrite}{%
\begin{algorithm}
\caption{Rewrite \textendash\ operator coordinating a cohort rewrite at height $h$}\label{alg:rewrite}
\pseudocode[linenumbering,codesize=\scriptsize]{%
\kw{collect}\ \text{the changes for one aligned subtree; } n_j \gets \text{new element, } \circ \text{ if unchanged} \\
\mathsf{tx}\text{: } \vin_0 = \text{tree output, plus funding for any increase} \\
\pcind \vout_0 = (S - \textstyle\sum o_j.\amt + \sum n_j.\amt,\ \TreeSpk(r', m)) \\
\pcind \vout_{1}, \dots \gets (o_j.\amt, \Ptr{o_j.\xk}) \text{ for each emptied leaf, in position order (a reclaimable leaf emptied under check 8 is not paid)} \\
\kw{for each}\ j \text{ whose leaf changes (not an amount-only increase):} \\
\pcind \kw{send}\ (\mathsf{tx}, o_j, n_j, \pi, r, m) \text{ to its owner} \\
\kw{each owner}\text{:}\ \kw{if}\ o_j \text{ is in } r \text{ (fold), and } n_j \text{ and every output are as agreed:} \\
\pcind \text{sign (and, under } x\text{, the exit key)} \\
\kw{if}\ \text{a required signature is missing:}\ \text{set } n_j \gets \circ,\ \text{rebuild, collect from line 2} \\
\sigma_p \gets \fn{Sign}_{\mathsf{OPERATOR\_KEY}}(\mathsf{tx});\ \text{assemble the witness;}\ \kw{broadcast}
}
\end{algorithm}%
The algorithm produces a transaction $\Phi_{\tr{REWRITE}}$ accepts. Lines 3 and 4
are checks 7 and 6: each emptied leaf is paid exactly, which is why the
operator cannot remove a leaf without paying it, and the output~0 amount
nets the subtree's old and new amounts. Lines 5 and 6 send the transaction only
to the owners whose leaves change in a way check~5 names, since a leaf
that only gains amount, unless it is $\mathsf{SEQ\_FINAL}$, and a leaf nobody
signs for pass through unchanged; positions in the reclaim form need the
operator alone (check~8). Lines 7 and 8 are the owner's side, which the spec
states for \tr{UPDATE}, whose signers MUST verify the full transaction
(spec~9.5), and which carries over to the cohort form by analogy; an
owner who emptied its position is paid by the covenant, and one who did
not sign keeps its leaf. Line~9 follows from the signature coverage: every
signature covers all outputs, so any change to the transaction
invalidates the others and the round restarts, as for a self-funded graft
(spec~9.2). The spec does not specify the message flow, when to give up
on an owner, or how requests are collected (design/batching.md describes
the transition, not the coordination).
}

% ------------------------------------------------------------------
\newcommand{\algEvict}{%
\begin{algorithm}
\caption{Evict \textendash\ operator, removing expired leaves}\label{alg:evict}
\pseudocode[linenumbering,codesize=\scriptsize]{%
\kw{admit}\ \text{a leaf to a cohort subtree only if } \ell.\ek = \mathsf{OPERATOR\_KEY}\text{,} \\
\pcind \text{and align the leases per cohort} \\
\kw{before}\ \ell.\ea\text{:}\ \text{notify the owner; offer an } \tr{UPDATE} \text{ extending } \ell.\ea \\
\kw{when}\ \lock \ge T = \max_j o_j.\ea \text{ over an aligned subtree:} \\
\pcind \mathsf{tx} \gets \Phi_{\tr{EVICT\_SUBTREE}} \text{ with the } 2^h \text{ elements published},\ \lock = T \\
\pcind \vout_{k_j} \gets (o_j.\amt, \Ptr{o_j.\xk}) \text{ for each leaf} \\
\pcind \vout_0 \gets (S - \textstyle\sum o_j.\amt, \TreeSpk(r', m));\ \text{sign with } \mathsf{OPERATOR\_KEY}\text{;} \\
\pcind \kw{broadcast},\ \text{paying the fee} \\
\kw{for}\ \text{a leaf that blocks the sweep (later lease, other evictor):} \\
\pcind \text{wait, or remove it with } \Phi_{\tr{EXIT}},\ e = 1
}
\end{algorithm}%
Lines 4 to 8 are the sweep of spec~7.7: one transaction pays every leaf
of the subtree in full to its exit key, so an eviction returns the
member's money and does not take it, at a small fraction of the per-leaf
cost of evicting each leaf alone (spec~9.4). The sweep requires a lock
time of at least the lease of every leaf in the subtree and
$\mathsf{OPERATOR\_KEY}$ as every evictor, which is why lines 1 and 2 make the
evictor key a condition of joining a cohort and align the leases: one
leaf with a later lease holds up its whole cohort, and one with another
evictor blocks it outright. Line~3 is the spec's advice (SHOULD) to notify
and to offer a renewal; a renewal by \tr{UPDATE} in place blocks the sweep
the same way, so the spec recommends relocating the renewed leaf into a
cohort whose lease matches, by an \tr{UPDATE} followed by a \tr{SPLIT} of
the whole balance or by a cooperative exit and a $\tr{JOIN}_0$
(spec~9.4). Lines 9 and 10 are the spec's remedy for a blocking leaf (spec~7.7). A
member whose lease is later than the cohort's can carve into the
cohort's empty positions and delay its sweep; the spec asks operators to
keep cohorts full or detach them, and gives no mechanism.
}

% ------------------------------------------------------------------
\newcommand{\algReceive}{%
\begin{algorithm}
\caption{Receive \textendash\ recipient $R$ and operator $O$, for one payment into a pair $(P, L)$}\label{alg:receive}
\pseudocode[linenumbering,codesize=\scriptsize]{%
R\text{:}\ \text{draw } \rho;\ \eta \gets \mathsf{SHA256}(\rho);\ \text{invoice on } \eta;\ \text{the HTLC } (\eta, a, E) \text{ reaches } O \\
O\text{:}\ \kw{require}\ E \ge T + \Dl + \text{time to settle from an on-chain preimage} \\
O\text{:}\ X \gets \text{settled level of } L + a \\
O\text{:}\ \kw{if}\ P.\amt < (X' - L.\amt) + \mathsf{DUST\_MIN} \text{ for } X' = X \text{ or an outstanding target:} \\
\pcind \text{top up } P \text{ by an amount-only increase} \\
O\text{:}\ \text{sign a } \tr{CLAIM}^{hm} \text{ voucher } V \text{ on } (\eta, X) \text{ against } P;\ \kw{send}\ V \text{ to } R \\
R\text{:}\ \kw{verify}\ P \text{ is reclaimable and } L\text{'s sibling, } L \text{ is } R\text{'s and not reclaimable,} \\
\pcind P.\xk = L.\ok = L.\xk \text{ is } R\text{'s key, and neither } P.\sq \text{ nor } L.\sq \text{ is } \mathsf{SEQ\_FINAL}, \\
\pcind V \text{ verifies under } P.\ok \text{ over } \mathsf{tree\_id},\ \hat P,\ \hat L \text{ and } X \text{ as in the confirmed tree}, \\
\pcind X = \text{settled level of } L + a,\ P.\amt + L.\amt - X' \ge \mathsf{DUST\_MIN} \text{ for every outstanding } X', \\
\pcind T \text{ at least } \Dl + \varepsilon \text{ away, and } L.\ea > T + \Dl\text{; else } \kw{refuse} \\
R\text{:}\ \text{release } \rho \text{ (one preimage at a time per pair); hand } (V, \rho) \text{ to its tower} \\
O\text{:}\ \kw{on}\ \text{learning } \rho\text{:}\ \text{settle the inbound HTLC; nothing reaches the chain}
}
\end{algorithm}%
\iffullversion
This is the construction of spec~9.6.3. The recipient holds a
\emph{receive pair}: a pending leaf $P$ at an even position, owned by the
operator, whose exit key is the recipient's, whose evictor is
$\mathsf{RECLAIM\_KEY}$ and whose deadline is $T$, and its sibling, the
landing leaf $L$, which the recipient owns and exits to. Line~1 has the
recipient draw the preimage, the party the HTLC pays, as the ordering
rules of \fn{Swap} require. Line~2 is rule~4 of the spec: $T$ must be
chosen so that the HTLC cannot expire before the operator learns a
preimage published late on-chain. Lines 3 to 5 are rule~1: the headroom is
required, since a voucher that drains $P$ exactly would be blocked by any
amount-only top-up, which leaves a sub-dust residual that
$\Phi_{\tr{CLAIM}}$ check~11 rejects, after which the operator could
reclaim the payment at $T$. Line~6 signs the merge voucher of
$\Phi_{\tr{CLAIM}}$ (checks~9 to 12): it raises $L$ to the absolute level
$X$ and lowers $P$ by the difference without advancing $P$'s $\sq$, so it
voids nothing and an older, lower voucher is a no-op after a newer one.
Lines 7 to 11 are the recipient's verification of rule~2, against the
\emph{confirmed} tree and before it releases anything, which also
excludes a pending leaf that is not $L$'s sibling or an $L$ that the
operator could evict inside the window; line~12 is the release rule, one
preimage at a time per pair, because the targets are levels of one leaf
and a second voucher that included an unsettled payment would hand the
recipient that payment without its preimage. Before line~12 the
operator has signed nothing it can lose, since $V$ needs a preimage that
only $R$ holds and publishing it on-chain hands it to $O$; after it $R$
can redeem alone. The spec fixes none of the invoice's expiry
arithmetic beyond rule~4 (design/pending-receive.md derives a minimum
final delta in terms of the window length).
\else
The recipient holds a \emph{receive pair}: a pending leaf $P$ (owned by
the operator, exit key the recipient's, evictor $\mathsf{RECLAIM\_KEY}$,
deadline $T$) and its sibling, the landing leaf $L$, which the recipient
owns (spec~9.6.3). Line~2 keeps the HTLC alive past a preimage published
late; lines 3 to 5 leave headroom, since a voucher that drains $P$
exactly is blocked by an amount-only top-up that leaves a sub-dust
residual. Line~6 signs a hash-locked merge voucher that raises $L$ to
the absolute level $X$ and lowers $P$ without advancing its $\sq$, so it
voids nothing and an older, lower voucher is a no-op. Lines 7 to 11 verify
against the \emph{confirmed} tree before anything is released, and line~12
releases one preimage at a time per pair. Before line~12 the operator has
signed nothing it can lose, since $V$ needs a preimage that only $R$
holds; after it $R$ can redeem alone.
\fi
}

% ------------------------------------------------------------------
\newcommand{\algSettle}{%
\begin{algorithm}
\caption{SettleWindow \textendash\ operator, at the close of a window of a pair $(P, L)$}\label{alg:settle}
\pseudocode[linenumbering,codesize=\scriptsize]{%
\kw{before}\ T - \Dl - \varepsilon\text{:}\ \kw{for each}\ \text{pair that settled amounts in the window:} \\
\pcind \text{raise } L \text{ to its settled level, either by redeeming the merge voucher } U \\
\pcind \text{that the operator signs itself,} \\
\pcind \text{or by a transfer in } \Phi_{\tr{REWRITE}} \text{ (reclaim form), batched across recipients} \\
\kw{hold}\ \text{every inbound HTLC with an outstanding merge voucher open} \\
\pcind \text{until the reclaim of } P \text{ has confirmed} \\
\kw{after}\ T\text{:}\ \text{run } \fn{Reclaim} \text{ on every } P \text{ of the window}
}
\end{algorithm}%
\iffullversion
Line~1 is the operator's duty under rule~3 of spec~9.6.3: before
$T - \Dl - \varepsilon$ it should raise $L$ to the level the recipient has settled in
the window, so that the recipient's vouchers move only what is still in
flight above it. The two ways of lines 2 to 4 are the voucher the operator signs
for itself, which needs no one else, and the transfer of the reclaim form
of the rewrite (spec~7.9 rule~8), which moves value from $P$ into the
sibling that the same key owns and exits to and needs the operator's
signature alone, so a cohort of recipients settles in one transaction
without any of them signing. Pairs that received nothing cost nothing.
The operator's promise is a SHOULD, because the recipient does not rely
on it: if $L$ has not reached the settled level by $T - \Dl - \varepsilon$, the
recipient or its tower MUST redeem the latest voucher (\fn{Watch}). Lines 5 and 6
are rule~4 and a MUST: until the reclaim confirms, the recipient can
still redeem and publish the preimage, and an HTLC already failed back
would leave the operator paying for a payment it never received. The
spec leaves open the length of a window, how many pairs a recipient
holds for payments across a deadline, and how the operator learns which
vouchers are outstanding.
\else
Lines 1 to 4 are rule~3 of spec~9.6.3: before $T - \Dl - \varepsilon$ the operator
raises $L$ to the settled level, by a voucher it signs itself or by the
transfer of the reclaim form of $\Phi_{\tr{REWRITE}}$, which needs the
operator's signature alone since it moves value only from $P$ into the
sibling that the same key owns and exits to; one transaction settles a
cohort of recipients. The recipient does not rely on it: if $L$ is short
at $T - \Dl - \varepsilon$, the recipient or a tower holding the latest voucher and the
preimage must redeem, the one liveness duty of the construction. Lines 5
and 6 are rule~4: until the reclaim of $P$ confirms, the recipient can
still redeem and publish the preimage, so an inbound HTLC with an
outstanding voucher stays open.
\fi
}

% ------------------------------------------------------------------
\newcommand{\algWatch}{%
\begin{algorithm}
\caption{Watch \textendash\ tower holding $(V, \rho)$ for a pair $(P, L)$}\label{alg:watch}
\pseudocode[linenumbering,codesize=\scriptsize]{%
\kw{on every}\ \text{spend of the tree output:}\ \fn{Follow}\ \text{the paths and the bytes of } P \text{ and } L \\
\kw{at}\ \text{height } T - \Dl - \varepsilon\text{ at the latest:} \\
\pcind \kw{if}\ L.\amt < \text{the settled level (confirmed)}\text{:} \\
\pcind\pcind \text{build } \Phi_{\tr{CLAIM}} \text{ (merge form) from the latest } V \text{ and } \rho \text{, own fee input} \\
\pcind\pcind \kw{broadcast};\ \text{if a conflicting spend confirms, rebuild as in } \fn{Exit} \\
\kw{stop}\ \text{once } L.\amt \ge X \text{ or the reclaim of } P \text{ has confirmed}
}
\end{algorithm}%
This is the one liveness duty of the recipient (spec~9.6.3 rule~3, which
puts the check on the recipient \emph{or} a tower holding the voucher and
preimage): once per window, at $T - \Dl - \varepsilon$ (line~2), confirm that $L$ holds
at least the settled level and, if it does not, redeem the latest voucher
with the preimage. A tower can do it because the merge voucher raises the
source's sibling, a leaf the recipient owns and exits to, so it moves
value only to a leaf the recipient alone can exit and needs no signature
beyond the voucher's (spec~7.5.4; design/towers.md); the tower holds no
key of the recipient's, and chooses only when to redeem and at what fee.
Line~1 is \fn{Follow}: the tower tracks the leaf bytes as well as the
paths, since a merge voucher binds $L$'s body, and a change to $L$ other
than its amount voids it (spec~7.5.4). Lines 3 to 5 are the redemption,
rebuilt against a new root if a conflicting spend wins, within the bound
$\Dl$ of spec~10.5. The spec gives no tower protocol (registration,
payment, how a voucher is handed over); design/towers.md lists it as not
designed, together with an incentive for towers.
}

% ------------------------------------------------------------------
\newcommand{\algReclaim}{%
\begin{algorithm}
\caption{Reclaim \textendash\ operator, taking back a pending leaf after its deadline}\label{alg:reclaim}
\pseudocode[linenumbering,codesize=\scriptsize]{%
\kw{when}\ \lock \ge P.\ea = T \text{ for a pending leaf } P \text{ with } P.\ek = \mathsf{RECLAIM\_KEY}\text{:} \\
\pcind \mathsf{tx} \gets \Phi_{\tr{RECLAIM}},\ \text{or the reclaim branch of } \Phi_{\tr{REWRITE}} \text{ replacing } P \text{ by the} \\
\pcind \text{next window's pending leaf} \\
\pcind \text{sign with } \mathsf{RECLAIM\_KEY} \text{ (} \Phi_{\tr{RECLAIM}}\text{) or } \mathsf{OPERATOR\_KEY} \text{ (the rewrite)};\ \kw{broadcast};\ \text{rebuild on a conflicting spend} \\
\kw{only after}\ \text{the reclaim has confirmed:}\ \text{fail upstream the HTLCs not settled}
}
\end{algorithm}%
Line~1 is the deadline of $\Phi_{\tr{RECLAIM}}$ (checks 2 and 3): a leaf
whose evictor is $\mathsf{RECLAIM\_KEY}$ is reclaimable, and after $T$ the
operator alone removes it with no payout, so that value in $P$ at its
deadline is the operator's, which is what makes the timeout of the
construction work. The batch form of lines 2 and 3 shares one transaction among
many pairs and writes the next window's pending leaf in the same
rewrite (spec~7.9 rule~8). A reclaimable leaf is never swept by
$\Phi_{\tr{EVICT\_SUBTREE}}$, whose evictor must be $\mathsf{OPERATOR\_KEY}$,
and an ordinary eviction of one still pays its exit key. Line~5 is rule~4
of spec~9.6.3: only after the reclaim confirms may the operator fail an
inbound HTLC whose merge voucher is outstanding, since until then the
recipient can still redeem and publish the preimage. The spec does not
say how the operator chooses between the two forms of line~2.
}
